\problemname{The Train Trip}
A group of people are going on a train trip and are wondering how they should sit.
$M$ pairs of people are friends and want to sit as close to each other as possible.
The train car they will sit in has $N$ rows and is designed as an $N \times 4$ grid, with 4 seats on each row.
There are $4N$ people in the group.
The people are numbered $1$ to $4N$.

Each pair of friends sitting at seats with Euclidean distance $L = \sqrt{dx^2 + dy^2}$
contributes $1/L^2$ to the group's total happiness.
Your task is to place the people so that the group's happiness is as large as possible.

\section*{Input}
The input consists of $10$ test cases.

The first line contains the number $T$ ($0 \le T \le 10$), which describes the test case number ($0$ for sample).
The second line contains the numbers $N$ and $M$ ($1 \le N \le 25\,000$, $1 \le M \le 100\,000$), the number of rows in the train car and the number of pairs of friends.
The following $M$ lines contain two integers $a$ and $b$ ($1 \le a,b \le 4N$, $a \ne b$), the numbers of two friends.

\section*{Output}
Print $N$ lines with 4 integers each.
Each line should contain the numbers of the four people sitting on that row.
All people must be placed somewhere.

\section*{Scoring}
Your program will be tested on $10$ test cases, and scored based on how it compares to \emph{all other contestants}.
If your submission on a given test case achieves a sum of happiness equal to $X$, and the best sum anyone has achieved on that test case is $Y$,
then your score on that test case is $10 \cdot (X / Y)^2$, rounded to two decimal places.

Your score for a submission is the sum of the scores for the test cases, and your total score is the score of your best submission.
It is thus not possible to combine test cases between different submissions; you must submit a single program that
maximizes your score on all test cases simultaneously.

During the contest, your score will be set relative to a judge solution;
after the contest ends, all solutions will be rejudged against the best contestant solutions.
If your solution is better than the one the judges put together, it is thus possible that you temporarily get more than $10$ points on a test case.
Your (temporary) total score will, however, be capped at $100$.

The sample test case will always give $0$ points.


\section*{Explanation of Sample}
In the sample, we want to place 8 people in a $2 \times 4$ grid.
The sample solution optimizes the distance for all friendships except 1 and 4 who sit at distance $\sqrt 3$ from each other.
The sum of happiness becomes $4 \cdot 1/1 + 1/3 \approx 4.33$.

A better solution can be achieved so that all pairs of friends sit at distance 1.
If that test case had been a real test case and another contestant had generated this solution
(total happiness $5$), then the test case would have given $10 \cdot (4.33 / 5)^2 \approx 7.51$ points.

\section*{Test Cases}
Here are descriptions of the $10$ test cases that appear:

\noindent
\begin{tabular}{| l | l | l |}
\hline
Test case & Constraints                       & Description \\ \hline
$1$       & $N = 10$, $M = 100$           & All edges are generated randomly. \\ \hline
$2$       & $N = 100$, $M = 500$          & All edges are generated randomly. \\ \hline
$3$       & $N = 100$, $M = 8\,000$       & All edges are generated randomly. \\ \hline
$4$       & $N = 8\,000$, $M \le 45\,000$ & $G$ is made up of groups of friends of size $1$-$5$ where everyone knows everyone. \\ \hline
$5$       & $N = 2\,500$, $M = 8\,000$    & $G$ contains no cycles. The maximum distance in $G$ is $2$. \\ \hline
$6$       & $N = 2\,500$, $M = 9\,800$    & $G$ contains no cycles. The maximum distance in $G$ is $2$. \\ \hline
$7$       & $N = 2\,500$, $M = 9\,999$    & $G$ contains no cycles. \\ \hline
$8$       & $N = 25\,000$, $M = 99\,999$  & $G$ contains no cycles. \\ \hline
$9$       & $N = 10\,000$, $M = 100\,000$ & All edges are generated randomly. \\ \hline
$10$      & $N = 25\,000$, $M = 100\,000$ & All edges are generated randomly. \\ \hline
\end{tabular}

By randomly we mean \href{https://en.wikipedia.org/wiki/Discrete_uniform_distribution}{uniformly at random}.
The graph of friendships is called $G$ in some of the descriptions.
